\section{Results from Earlier Builds}
\label{sec:archive-results}
These archived results retain their original builds, storage modes,
and measurement endpoints. The main paper reports continuous-payment
and representative P results on \texttt{f856c7b}; historical-reader
results retain \texttt{721800c}, and the injected-latency study retains
\texttt{9879f64}. Results from different builds and workloads are not
combined into a capacity estimate for one configuration.

\subsection{Continuous Spending Without Per-Hop Settlement}
\label{archive:sec:eval-continuation}

We compare certified continuation with waiting for public execution between
hops. On the baseline build, a
100-hop chain reaches its last certified receipt in a median 158.846\,ms,
compared with 65.204\,s when each hop waits for public confirmation. Three
wallets transfer 100 CAL in rotation, constructing each payment online from
its predecessor's output and paying fees from their own final FUEL. Three
repetitions per mode give six runs and 600 payments on the nine-service
in-memory configuration. Both modes end at the last certified receipt, and
the waiting mode adds a public confirmation \emph{between} hops. With
committee commit, flush, and gossip settings of 250, 10, and 10\,ms, public
execution and its observation cost about 0.65\,s per hop. Certified
continuation removes that wait from the recipient's path.

Figure~\ref{archive:fig:continuation} shows the six whole-chain durations. The fast
runs take 158.492--159.905\,ms. All 297 fast-mode successor requests consume
the preceding certified output and are sent after its certified receipt but
before its public completion is observed. A one-input, one-output TXCer is
779 bytes, and requests are 1,979 bytes at the first hop and 2,766 bytes
thereafter. For this transaction shape, the direct evidence stays the same
size across all 100 hops, because each request carries the certificates of its
direct inputs rather than the ancestor chain.

\begin{figure}[t]
  \centering
  \includegraphics[width=\columnwidth]{continuation-current.pdf}
  \caption{Time through the last certified receipt of a 100-hop chain on
    the baseline build. Dots are three independent runs per mode;
    horizontal bars mark their medians. The logarithmic time axis includes
    the same endpoint in both modes. Only the waiting baseline requires
    public confirmation between successive payments.}
  \label{archive:fig:continuation}
\end{figure}

To test whether the advantage depends on in-memory member storage, three paired 100-hop trials on a nine-service disk-backed
configuration with NoSync wallets and user-paid FUEL differ only in member
storage (NoSync or synchronous; order N/S, S/N, N/S), and each is followed by
a synchronous waiting chain. Median durations are 681.127\,ms for NoSync
continuation, 4.471\,s for synchronous continuation, and 65.013\,s for
synchronous waiting. Certified continuation therefore remains 14.5 times
faster than waiting when member commits are synchronous. All 900 payments
complete with agreeing committee states.

On an archived delivery-gate build, a gate after certification waits for three
members to INSTALL complete material or for authenticated public success.
The gate collects no second signature quorum, and delivering immediately
instead shortens continuation. At 200
payments/s for 30\,s, three paired runs per mode each complete 6,000
payments, and the median of run-level receipt P50s falls from 2.167\,ms with
the gate to 1.300\,ms (40.0\%). Separate 100-hop trials fall from 287.983 to
200.356\,ms (30.43\%). Authorization is unchanged.
Fig.~\ref{fig:delivery-gate} shows the run-level results;
Section~\ref{sec:supp-methods} gives the methods.

\subsection{Funded Closure and Repayment}
\label{archive:sec:eval-liabilities}

This experiment checks that compensation leaves successor payments in place
and that each compensation is repaid exactly once when its source later
executes.

\noindent\textbf{Compensation and repayment matrix (baseline build).}
Nine runs use seeds 23, 37, and 59 at withheld-parent rates of 0\%, 1\%, and
5\%. Each run offers 20 two-hop units/s for 60\,s (1,200 units, 2,400
payments) against a fixed 60,000-CAL reserve with user-paid FUEL and no
top-up. The nine-service deployment uses disk-backed committee application
state and BlockStore, in-memory members and gateway, and NoSync wallets. The
driver releases each withheld parent only after its compensation command is committed
and its historical revision is materialized on all four replicas, so every withheld source is compensated
after the 30-s public-time deadline and executes afterwards.

All 21,600 payments close. The 216 compensations (12 per 1\% run, 60 per 5\%
run) are each followed by repayment. At the end of every run no obligation is
Open or Repaired, the gap is zero, net Spent and Reserved are zero, and the
reserve is back at 60,000\,CAL (Table~\ref{archive:tab:eval-audit}). The lowest
sampled reserve is 59,800--59,900\,CAL at 1\% and 59,600\,CAL at 5\%. It
reflects the driver's release schedule rather than a reserve-sizing result,
and Fig.~\ref{fig:reserve-recovery} shows the three 5\% trajectories. Run-level receipt P50s
for normal child payments are 1.185--1.251\,ms. Measured from the deadline, the
committee commit, which in this build carries the representation, is observed
at a median of 1.003--1.254\,s, and physical installation of the revised
body on all four replicas at a median of 1.257--1.761\,s (P95 up to 2.015\,s); both include 250-ms observation
polling. Each normal payment spends 94 FUEL and each compensation adds 5, and
rewards, burn, and refunds sum to the fee caps with zero Held.

\begin{table}[t]
  \caption{Closing audit of the compensation and repayment matrix. Each
    column is reproduced by all three seeds; each run has 2,400 payments.
    Reserve values are CAL, fee values are in thousands of FUEL, and fee
    caps sum the per-payment maxima.}
  \label{archive:tab:eval-audit}
  \centering
  \small
  \begin{tabular}{@{}lrrr@{}}
    \toprule
    Withheld parents & 0\% & 1\% & 5\% \\
    \midrule
    Compensations / repayments & 0 & 12 & 60 \\
    Paid / repaid CAL & 0 & 1,200 & 6,000 \\
    Final reserve & 60,000 & 60,000 & 60,000 \\
    Unclosed / net Spent & 0 & 0 & 0 \\
    Total fee caps & 2,400 & 2,400 & 2,400 \\
    Rewards & 201.600 & 201.660 & 201.900 \\
    Burned & 24.000 & 24.000 & 24.000 \\
    Refunded & 2,174.400 & 2,174.340 & 2,174.100 \\
    Remaining Held & 0 & 0 & 0 \\
    \bottomrule
  \end{tabular}
\end{table}

Ledger-level regression covers arrival orders directly: all 24 arrival orders
of a four-hop chain, with and without due compensation, under three issuer
sequences give 144 trajectories, and 48 split--merge trajectories vary all
arrival orders of a 40/60 split, two branches, and their merge under two
compensation modes. Every step in the 192 trajectories conserves the
independently computed asset projection. Every trajectory ends with each
issuer's reserve restored, Spent and Reserved zero, $\mathsf{Paid}=\mathsf{Rec}$,
and the consumed inputs still consumed.

Two fourteen-service controls separate permanent loss from delayed
repayment. When the parent never arrives, the grandchild keeps its 100 CAL
and the responsible issuer's reserve stays 100 CAL lower. When a
cross-organization chain arrives grandchild first, the intermediate source
closes its own direct obligation, and the root issuer's compensated 100 CAL
is repaid when the root later executes. The four committee states agree in
both cases, and an absent parent's own fee and work reservations can remain
outstanding.

Three paired
fourteen-service runs on the initial baseline build, before the batch-result
tag separation (4,000 payments at 100 payments/s each), end with identical
owner holdings and reserve balances in normal and withheld-source runs,
consistent with the CAL delay-neutrality theorem in the main paper. Concurrent compensation
raises receipt P95 by a factor of 1.98--2.21, and the pairs do not isolate
the responsible stage.

A storage control with NoSync wallets isolates member
commits: receipt P50 is 1.0\,ms with NoSync members and 43--52\,ms with
synchronous commits.

An earlier-build capital series with one run per
condition and no compensation shows that external additions of
10,800--19,700\,CAL sustain 300-s runs of 6,000 two-hop units from a
3,600-CAL start, whereas a fixed 3,600-CAL control closes 46 of 6,000 units.

\subsection{Continuation Across Registered Organizations}
\label{archive:sec:eval-network}

To measure the effect of communication distance, we use a two-organization
suite on an archived build recorded in the supplement. Its 24 runs complete
342,000 payments, with two wallet processes holding 32
identities each across 64 lanes and each organization supplying 50
payments/s. The latency comparison uses ten-hop same- and
cross-organization workloads at 0, 20, and 100\,ms injected HTTP RTT, three
repetitions per condition. Each proxy direction adds half the RTT, and
intra-organization certification and committee P2P receive no delay
injection.

\begin{figure}[t]
  \centering
  \includegraphics[width=\columnwidth]{network.pdf}
  \caption{ACK latency (ms) versus injected cross-site HTTP RTT (ms).
    Each point is the median of three run-level P50s for the first
    120-s cohort; error bars span those P50s' minimum and maximum.
    Aggregate offered rate is 100 payments/s. Committee P2P is undelayed.}
  \label{archive:fig:network}
\end{figure}

For the common first-120-s cohort, the median of three run-level cross-organization
ACK P50s is 1.632, 21.658, and 101.967\,ms, respectively, and same-organization
values span 1.615--1.744\,ms (Fig.~\ref{archive:fig:network}). Cross-organization ACK P50
tracks the injected RTT, with a measured excess of 1.6--2.0\,ms. An ACK includes the
return path, but a recipient may start its successor on acceptance without
waiting for the ACK to reach the previous sender. At 100-ms injected RTT in the cross-organization workload,
three 300-s runs each complete 30,000 payments; pending peaks are 106--110 and
drain without sustained tail growth. The experiment isolates HTTP path delay
within one physical host.

\subsection{Operating Capacity and Service Availability}
\label{archive:sec:eval-operating}

\noindent\textbf{Independent-input throughput (baseline build).}
Three 30,000-payment runs target 2,000 payments/s on the nine-service
in-memory configuration with NoSync wallets and authorized organization-paid
FUEL. Member processes use a 200\% Go GC target and 16 runtime processors. All
payments complete and the four committee states agree. Completion rates
including drain are 1,750.27, 1,820.91, and 1,887.31 payments/s over
15.9--17.1\,s, with receipt P50s of 2.646, 2.842, and 2.824\,ms and P95s of
47.9, 49.1, and 48.0\,ms. The driver recorded 6,551--7,570 payments that had to
wait for a fast-receipt admission slot, and the P95 dispatch lag was
179--229\,ms, so the offered load was not uniform at 2,000 payments/s. These
are short-run completion measurements on the baseline build. They are not
remeasured for the decision-first build, whose cost on ordinary payments is
tested by the six-run comparison above. The supplement identifies the build
and duration of the 150,000-payment series.

On archived builds, redundant certification and evidence replicated by INSTALL
sustain service under the specified failures. Nine runs at 200 payments/s complete
216,000 payments with one member delayed or paused for 30\,s. Gateway-loss
controls in which one member holds a complete copy replicated by INSTALL
show backup submission 2.503--2.508\,s after INSTALL.
Conflict controls produce neither dual conflicting QCs nor repeated
principal/fee effects, while partial approvals remain a capacity cost.

The multi-identity measurements in Section~\ref{sec:supp-methods} report
per-payment CPU and HTTP-payload costs; 2,048 identities share clients.
Historical memory growth is reported separately from backlog. The same
section gives an LND regtest reference, whose completion event,
persistence, and workload concurrency differ from the Next measurements.
